% Options for packages loaded elsewhere
\PassOptionsToPackage{unicode}{hyperref}
\PassOptionsToPackage{hyphens}{url}
%
\documentclass[
]{article}
\usepackage{amsmath,amssymb}
\usepackage{setspace}
\usepackage{iftex}
\ifPDFTeX
  \usepackage[T1]{fontenc}
  \usepackage[utf8]{inputenc}
  \usepackage{textcomp} % provide euro and other symbols
\else % if luatex or xetex
  \usepackage{unicode-math} % this also loads fontspec
  \defaultfontfeatures{Scale=MatchLowercase}
  \defaultfontfeatures[\rmfamily]{Ligatures=TeX,Scale=1}
\fi
\usepackage{lmodern}
\ifPDFTeX\else
  % xetex/luatex font selection
\fi
% Use upquote if available, for straight quotes in verbatim environments
\IfFileExists{upquote.sty}{\usepackage{upquote}}{}
\IfFileExists{microtype.sty}{% use microtype if available
  \usepackage[]{microtype}
  \UseMicrotypeSet[protrusion]{basicmath} % disable protrusion for tt fonts
}{}
\makeatletter
\@ifundefined{KOMAClassName}{% if non-KOMA class
  \IfFileExists{parskip.sty}{%
    \usepackage{parskip}
  }{% else
    \setlength{\parindent}{0pt}
    \setlength{\parskip}{6pt plus 2pt minus 1pt}}
}{% if KOMA class
  \KOMAoptions{parskip=half}}
\makeatother
\usepackage{xcolor}
\usepackage[margin=1in]{geometry}
\usepackage{longtable,booktabs,array}
\usepackage{calc} % for calculating minipage widths
% Correct order of tables after \paragraph or \subparagraph
\usepackage{etoolbox}
\makeatletter
\patchcmd\longtable{\par}{\if@noskipsec\mbox{}\fi\par}{}{}
\makeatother
% Allow footnotes in longtable head/foot
\usepackage{graphicx}
\makeatletter
\def\maxwidth{\ifdim\Gin@nat@width>\linewidth\linewidth\else\Gin@nat@width\fi}
\def\maxheight{\ifdim\Gin@nat@height>\textheight\textheight\else\Gin@nat@height\fi}
\makeatother
% Scale images if necessary, so that they will not overflow the page
% margins by default, and it is still possible to overwrite the defaults
% using explicit options in \includegraphics[width, height, ...]{}
\setkeys{Gin}{width=\maxwidth,height=\maxheight,keepaspectratio}
% Set default figure placement to htbp
\makeatletter
\def\fps@figure{htbp}
\makeatother
\setlength{\emergencystretch}{3em} % prevent overfull lines
\providecommand{\tightlist}{%
  \setlength{\itemsep}{0pt}\setlength{\parskip}{0pt}}
\setcounter{secnumdepth}{-\maxdimen} % remove section numbering
\newlength{\cslhangindent}
\setlength{\cslhangindent}{1.5em}
\newlength{\csllabelwidth}
\setlength{\csllabelwidth}{3em}
\newlength{\cslentryspacingunit} % times entry-spacing
\setlength{\cslentryspacingunit}{\parskip}
\newenvironment{CSLReferences}[2] % #1 hanging-ident, #2 entry spacing
 {% don't indent paragraphs
  \setlength{\parindent}{0pt}
  % turn on hanging indent if param 1 is 1
  \ifodd #1
  \let\oldpar\par
  \def\par{\hangindent=\cslhangindent\oldpar}
  \fi
  % set entry spacing
  \setlength{\parskip}{#2\cslentryspacingunit}
 }%
 {}
\usepackage{calc}
\newcommand{\CSLBlock}[1]{#1\hfill\break}
\newcommand{\CSLLeftMargin}[1]{\parbox[t]{\csllabelwidth}{#1}}
\newcommand{\CSLRightInline}[1]{\parbox[t]{\linewidth - \csllabelwidth}{#1}\break}
\newcommand{\CSLIndent}[1]{\hspace{\cslhangindent}#1}
\usepackage[left]{lineno}
\linenumbers
\usepackage{hyperref}
\hypersetup{backref=true}
\usepackage{setspace}\doublespacing
\usepackage{booktabs}
\usepackage{longtable}
\usepackage{array}
\usepackage{multirow}
\usepackage{wrapfig}
\usepackage{float}
\usepackage{colortbl}
\usepackage{pdflscape}
\usepackage{tabu}
\usepackage{threeparttable}
\usepackage{threeparttablex}
\usepackage[normalem]{ulem}
\usepackage{makecell}
\usepackage{xcolor}
\IfFileExists{bookmark.sty}{\usepackage{bookmark}}{\usepackage{hyperref}}
\urlstyle{same}
\hypersetup{
  pdftitle={Response diversity is a major driver of temporal stability in complex food webs},
  pdfauthor={Alain Danet (1,2), Sonia Kéfi (3,4), Thomas F. Johnson (1), Andrew P. Beckerman (1)},
  hidelinks,
  pdfcreator={LaTeX via pandoc}}

\title{Response diversity is a major driver of temporal stability in complex food webs}
\author{Alain Danet (1,2), Sonia Kéfi (3,4), Thomas F. Johnson (1), Andrew P. Beckerman (1)}
\date{}

\begin{document}
\maketitle

\setstretch{2}
\begin{singlespace}

\begin{itemize}
\item
  Institutional adress:

  \begin{itemize}
  \tightlist
  \item
    1 School of Biosciences, Ecology and Evolutionary Biology, University of Sheffield, Western Bank, Sheffield, S10 2TN, UK
  \item
    2 DECOD (Ecosystem Dynamics and Sustainability), INRAE, Institut Agro, IFREMER, Rennes, France
  \item
    3 ISEM, CNRS, Univ. Montpellier, IRD, EPHE, Montpellier, France
  \item
    4 Santa Fe Institute, 1399 Hyde Park Road, Santa Fe, NM 87501, USA
  \end{itemize}
\item
  Email address:

  \begin{itemize}
  \tightlist
  \item
    Alain Danet (corresponding author): \href{mailto:alain.danet@inrae.fr}{\nolinkurl{alain.danet@inrae.fr}}
  \item
    Sonia Kéfi: \href{mailto:sonia.kefi@umontpellier.fr}{\nolinkurl{sonia.kefi@umontpellier.fr}}
  \item
    Thomas F Johnson: \href{mailto:t.f.johnson@sheffield.ac.uk}{\nolinkurl{t.f.johnson@sheffield.ac.uk}}
  \item
    Andrew P Beckerman: \href{mailto:a.beckerman@sheffield.ac.uk}{\nolinkurl{a.beckerman@sheffield.ac.uk}}
  \end{itemize}
\item
  Short running title: Response diversity drives food web stability
\item
  Keywords: food webs, temporal stability, theoretical modelling, response
  diversity
\item
  Type of article: Research
\item
  Number of words: Abstract: 177, Main text: 6965
\item
  Number of references: 68, Number of figures: 2, Number of tables: 0
\item
  Corresponding author: Alain Danet, School of Biosciences, Ecology and
  Evolutionary Biology, University of Sheffield, Western Bank, Sheffield, S10
  2TN, UK; e-mail address: \href{mailto:alain.danet@inrae.fr}{\nolinkurl{alain.danet@inrae.fr}}
\end{itemize}

\end{singlespace}

\pagebreak

\hypertarget{abstract}{%
\section{Abstract}\label{abstract}}

Global change constitutes a major threat to biodiversity and ecosystem
functions, and places the temporal stability of ecological communities at risk.
Classic theory identifies species richness and food web structure as key
drivers of temporal stability, while recent work highlights response
diversity---variation in species' responses to environmental perturbations---as a
critical stabilizer via asynchrony in population fluctuations. However, how
these mechanisms interact in complex, multi-trophic communities remains
unresolved. Using a stochastic, bioenergenetic food web model, we integrate
these multiple bodies of theory to reveal that response diversity is a major
driver of community stability. Moreover, our integrated theory reveals that
positive stability-richness relationships emerge only in the presence of
response diversity. In contrast to previous work, we also find that food web
structure is only a secondary driver of community stability, but interacts with
response diversity to determine the sign of the stability-richness
relationship. Our study reveals identifiable pathways by which food web
structure and response diversity drive community stability and raises concerns
about how the loss of response diversity may lead to a breakdown of community
stability.

\pagebreak

\hypertarget{introduction}{%
\section{Introduction}\label{introduction}}

Understanding how ecological communities withstand environmental perturbations
and thus remain stable is a central goal in ecology {[}\protect\hyperlink{ref-macarthur_fluctuations_1955}{1},\protect\hyperlink{ref-may_will_1972}{2}{]}. Theory and experiments show that species diversity enhances
temporal stability by buffering perturbations {[}\protect\hyperlink{ref-yachi_biodiversity_1999}{3}--\protect\hyperlink{ref-thebault_stability_2010}{11}{]}. This stability arises from two components: (1)
population-level stability, which measures the temporal variability in biomass
averaged over all species in the community and (2) asynchrony, which captures the
degree to which species' fluctuations offset one another
{[}\protect\hyperlink{ref-zhao_biodiversity_2022}{8},\protect\hyperlink{ref-thibaut_understanding_2013}{12}{]}. Asynchrony underpins
diversity's stabilizing effects {[}\protect\hyperlink{ref-yachi_biodiversity_1999}{3},\protect\hyperlink{ref-de_mazancourt_predicting_2013}{6}--\protect\hyperlink{ref-zhao_biodiversity_2022}{8},\protect\hyperlink{ref-thibaut_understanding_2013}{12},\protect\hyperlink{ref-olivier_urbanization_2020}{13}{]}, and is
hypothetised to be mainly driven by response diversity, the diversity of
species' responses to environmental fluctuations {[}\protect\hyperlink{ref-elmqvist_response_2003}{14}--\protect\hyperlink{ref-polazzo_measuring_2024}{17}{]}.

Our understanding of the mechanisms driving temporal stability and
stability-richness relationships remains largely restricted to simplified
assemblages such as competitive {[}\protect\hyperlink{ref-yachi_biodiversity_1999}{3}--\protect\hyperlink{ref-zhao_biodiversity_2022}{8}{]} and plant-herbivore systems
{[}\protect\hyperlink{ref-downing_ecosystem_2002}{9}--\protect\hyperlink{ref-thebault_stability_2010}{11}{]}. This is problematic as complex assemblages like food
webs may behave differently as trophic interactions create strong
interdependence in species' dynamics {[}\protect\hyperlink{ref-raimondo_interspecific_2004}{18}--\protect\hyperlink{ref-fahimipour_compensation_2017}{22}{]}. Multiple lines of evidence support that trophic
interactions can alter community stability {[}\protect\hyperlink{ref-brose_allometric_2006}{23}--\protect\hyperlink{ref-eschenbrenner_diversity_2023}{26}{]}
through population stability {[}\protect\hyperlink{ref-may_will_1972}{2},\protect\hyperlink{ref-mccann_weak_1998}{20},\protect\hyperlink{ref-mccann_reevaluating_1997}{27}--\protect\hyperlink{ref-thebault_trophic_2005}{29}{]} and
asynchrony, and that food web structure can interact with response diversity
and species richness, the main drivers of temporal stability
{[}\protect\hyperlink{ref-mccann_weak_1998}{20}--\protect\hyperlink{ref-fahimipour_compensation_2017}{22}{]}. In the following paragraphs, we review the
evidence around this array of theory and then present our in-silico experiment
aiming to understand how these mechanisms combine to simultaneously affect
temporal stability of complex, multitrophic, ecological communities.

Stability at the population level is a key component of community stability.
Previous theoretical and empirical studies highlighted that population
stability decreases strongly with environmental stochasticity, whilst species
richness can have mixed effects, where a rise in richness has been observed to
both decrease and increase population stability {[}\protect\hyperlink{ref-tilman_biodiversity_2006}{30}--\protect\hyperlink{ref-houlahan_negative_2018}{32}{]}. Food web structure can
also mediate these effects as higher trophic levels may stabilize populations
via lower mortality {[}\protect\hyperlink{ref-brose_allometric_2006}{23},\protect\hyperlink{ref-hatton_linking_2019}{33}{]} or top-down
control {[}\protect\hyperlink{ref-shanafelt_stability_2018}{24}{]}, while strong interactions and high
connectance often decrease population stability {[}\protect\hyperlink{ref-may_will_1972}{2},\protect\hyperlink{ref-mccann_weak_1998}{20},\protect\hyperlink{ref-mccann_reevaluating_1997}{27},\protect\hyperlink{ref-thebault_trophic_2005}{29}{]}.
Crucially, these population responses to community structure are further
mediated by asynchrony to shape community stability
{[}\protect\hyperlink{ref-zhao_biodiversity_2022}{8},\protect\hyperlink{ref-thibaut_understanding_2013}{12}{]}.

This asynchrony emerges from two pathways in the absence of strong demographic
stochasticity {[}\protect\hyperlink{ref-loreau_biodiversity_2013}{31},\protect\hyperlink{ref-loreau_species_2008}{34}{]}: (i) response
diversity to environmental changes and (ii) species interactions. Response
diversity, arising from differences in species' niches and environmental
preferences, leads to varied reactions to the same environmental perturbations
{[}\protect\hyperlink{ref-elmqvist_response_2003}{14},\protect\hyperlink{ref-ross_how_2023}{16}{]}. For example, within a community,
some species may increase in abundance under warmer conditions while others
decline, depending on their thermal tolerances {[}\protect\hyperlink{ref-frund_response_2013}{35}{]}. Using a
trait-based or non trait-based approach, previous studies showed that response
diversity drives asynchrony {[}\protect\hyperlink{ref-thibaut_diversity_2012}{7},\protect\hyperlink{ref-thibaut_understanding_2013}{12},\protect\hyperlink{ref-sasaki_species_2019}{15}{]}, and thereby stability
{[}\protect\hyperlink{ref-yachi_biodiversity_1999}{3},\protect\hyperlink{ref-thibaut_diversity_2012}{7},\protect\hyperlink{ref-thibaut_understanding_2013}{12}{]}.

Response diversity modulates asynchrony through two distinct mechanisms:
portfolio effects and compensatory dynamics. Portfolio effects occur when
species exhibit independent responses to environmental fluctuations, leading to
statistical averaging that stabilizes aggregate community biomass, thereby
increasing community stability {[}\protect\hyperlink{ref-doak_statistical_1998}{4}{]}. This effect represents
statistical averaging weighted by response diversity, and is strongest when
species' environmental responses are uncorrelated and weakest when response
diversity is low (i.e., when species' responses are positively correlated).
Compensatory dynamics, by contrast, arise when species show negatively
correlated responses, providing additional stabilization
{[}\protect\hyperlink{ref-zhao_biodiversity_2022}{8}{]}. Empirical work in grasslands has demonstrated that
portfolio effects---driven primarily by response diversity---constitute the
dominant mechanism behind diversity's stabilizing effects
{[}\protect\hyperlink{ref-doak_statistical_1998}{4},\protect\hyperlink{ref-zhao_biodiversity_2022}{8}{]}.

Originally defined as the diversity of species responses within functional
groups {[}\protect\hyperlink{ref-elmqvist_response_2003}{14}{]}, the role of response diversity on asynchrony
via portfolio effects and compensatory dynamics in multi-trophic systems
remains poorly understood. Theory suggests that trophic interactions might
alter the effects of response diversity on asynchrony, either synchronizing
prey dynamics through top-down control {[}\protect\hyperlink{ref-raimondo_interspecific_2004}{18},\protect\hyperlink{ref-korpimaki_predatorinduced_2005}{19}{]} or enhancing compensatory dynamics via prey
switching and trophic cascades {[}\protect\hyperlink{ref-mccann_weak_1998}{20}--\protect\hyperlink{ref-fahimipour_compensation_2017}{22}{]}. These food web
structures might influence stability-richness relationships, where for example,
species-rich communities often contain more trophic levels. This makes it
challenging to disentangle the effects of richness from those of food web
structure {[}\protect\hyperlink{ref-danet_species_2021}{25},\protect\hyperlink{ref-eschenbrenner_diversity_2023}{26},\protect\hyperlink{ref-thebault_trophic_2005}{29}{]} per se {[}\protect\hyperlink{ref-dunne_network_2006}{36},\protect\hyperlink{ref-galiana_spatial_2018}{37}{]}.

To reconcile this wide range of theory defined across multiple ecological
scales and decompose the relative contribution of each process to stability and
stability-richness relationships, we here integrate these sets of processes
into an extended bioenergetic food web model {[}\protect\hyperlink{ref-brose_allometric_2006}{23},\protect\hyperlink{ref-yodzis_body_1992}{38}{]} that incorporates
environmental stochasticity and response diversity {[}\protect\hyperlink{ref-vasseur_environmental_2007}{39}--\protect\hyperlink{ref-ripa_food_2003}{41}{]}. We investigate
relationships among multiple processes across 45,000 in-silico communities with
varying levels of species richness, trophic complexity, connectance, and where
we experimentally manipulate interaction strengths, response diversity, and
environmental stochasticity. This in-silico experiment allowed us to
specifically evaluate whether response diversity, species richness, and food
web structure act additively or interactively to drive community stability, and
under what conditions these factors lead to positive or negative
stability-richness relationships. Finally, using a Structural Equation Model
and linear models, we test how the factors linked to environmental fluctuations
and their buffering (environmental stochasticity and response diversity) and
the ones mediating them (species richness and food web structure) affect the
components of community stability (population stability and asynchrony).

\hypertarget{methods}{%
\section{Methods}\label{methods}}

\hypertarget{bioenergenetic-model}{%
\subsection{Bioenergenetic model}\label{bioenergenetic-model}}

We simulated the dynamics of complex food webs using the allometric,
bioenergetic model {[}\protect\hyperlink{ref-brose_allometric_2006}{23},\protect\hyperlink{ref-yodzis_body_1992}{38},\protect\hyperlink{ref-rooney_homage_2006}{42},\protect\hyperlink{ref-lajaaiti_ecologicalnetworksdynamicsjl_2025}{43}{]}. This model is widely used to simulate the
dynamics of complex ecological communities because of the simplicity of its
parametrisation, where metabolism, growth and foraging rates are based on
species body mass and metabolic types using the metabolic theory of ecology
{[}\protect\hyperlink{ref-brown_toward_2004}{44}{]}. The model describes species biomass (\(B_i\)) dynamics over
time of the primary producers (\(B_{i \in \{\mathrm{prod.}\}}\), eq.
\eqref{eq:prod}) and the consumers (\(B_{i \in \{\mathrm{cons.}\}}\),
eq.\eqref{eq:cons}):

\begin{equation}
    \label{eq:cons}
        \frac{dB_{i \in \{\mathrm{cons.}\}}}{dt} =
        \sum_{j \in \{\mathrm{res.}\}_i} x_i y_i B_i F_{ij}
        - \sum_{j \in \{\mathrm{cons.}\}_i} \frac{x_j y_j B_j F_{ji}}{e_{ij}}
        - x_i B_i - d_i B_i
\end{equation}

\begin{equation}
    \label{eq:prod}
    \frac{dB_{i \in \{\mathrm{prod.}\}}}{dt} =
    r_i B_i G_i
    - \sum_{j \in \{\mathrm{cons.}\}_i} \frac{x_j y_j B_j F_{ji}}{e_{ij}}
    - d_i B_i
\end{equation}

The consumers gain biomass by consuming their resources, where \(j\in\{\mathrm{res.}\}_i\)
are the resources of the consumer \(i\), at a rate that depends on their
metabolic rate (\(x_i\)), maximum consumption rate (\(y_i\)) and on the functional
response (\(F_{ij}\)) describing how the rate of consumption of a consumer \(i\) on
one of its resources \(j\) vary with the biomass of this resource. The consumer
metabolic rate was derived from ectotherm mass specific metabolic rate (\(a_x = 0.88\), Table S1) and species body masses (\(M_i\)): \(x_i=a_xM_i^b\), with \(b=-1/4\)
as defined in the allometric theory of metabolism {[}\protect\hyperlink{ref-brose_allometric_2006}{23},\protect\hyperlink{ref-brown_toward_2004}{44}{]}.

Species then lose biomass by being consumed, where \(j \in \{\mathrm{cons.}\}_i\) are the consumers
of the species \(i\), \(e_{ij}\) is the assimilation efficiency of the consumer \(j\) on the
resource \(i\) {[}\protect\hyperlink{ref-brose_allometric_2006}{23},\protect\hyperlink{ref-rooney_homage_2006}{42},\protect\hyperlink{ref-lajaaiti_ecologicalnetworksdynamicsjl_2025}{43}{]}. Consumers also lose biomass over
time through metabolic losses (\(-x_iB_i\)) and natural mortality (\(-d_iB_i\)). The
primary producers gain biomass over time with a growth rate (\(r_i\), eq. \eqref{eq:prod}) and a
functional response (\(G_i\)) describing how the growth of the producers vary with
their biomass. They lose biomass by being consumed by consumers and by natural
mortality.

The growth of the primary producers (eq. \eqref{eq:growth}) is logistic,

\begin{equation}
    \label{eq:growth}
    G_i = (1 - \dfrac{\sum_{j \in \{\mathrm{prod.}\}} \alpha_{ij} B_j}{K_i})
\end{equation}

where the growth rate is maximal when \(\sum_{j \in \{\mathrm{prod.}\}} \alpha_{ij} B_j\) approaches \(0\) and null when it approaches \(K_i\). \(\alpha_{ij}\) is the per
capita effect of the producer \(j\) on the producer \(i\), \(K_i\) is the carrying
capacity for the producer \(j\). The feeding rates of a consumer feeding on
a resource (eq. \eqref{eq:holling}) is described by a functional response,

\begin{equation}
    \label{eq:holling}
    F_{ij} = \frac{\omega_{ij} B_j^h}
    {B_0^h + c_i B_i
        + \sum_{k \in \{\mathrm{res.}\}_i} \omega_{ik} B_k^h}
\end{equation}

which depends on the relative preference of the consumer on the resource
(\(\omega_{ij}\)) is limited by a half-saturation rate (\(B_0\)), intraspecific
interference coefficient (\(c_i\)), and on the availability of its other resources
(\(\sum_{k \in \{\mathrm{res.}\}_i} \omega_{ik} B_k^h\)). Finally, the \(h\)
exponent controls the shape of the functional response from a saturating
function (\(h = 1\), Holling type II) to a sigmoid (\(h = 2\), Holling type III).

\hypertarget{environmental-stochasticity-and-response-diversity}{%
\subsection{Environmental stochasticity and response diversity}\label{environmental-stochasticity-and-response-diversity}}

We added a stochastic natural mortality rate to the species dynamics,
representing a fluctuating environment. Based on the literature and previous
models {[}\protect\hyperlink{ref-hatton_linking_2019}{33},\protect\hyperlink{ref-vasseur_environmental_2007}{39}{]}, we assumed that natural mortality rates
scales inversely with the species body mass (\(d_i = d_0M_i^{-1/4}\)), as all
other physiological parameters. Following previous work, we used a basal
natural mortality rate of \(d_0 = 0.4\) {[}\protect\hyperlink{ref-vasseur_environmental_2007}{39}{]}.

Introducing environmental stochasticity via mortality rates allows to consider
stochasticity without modifying species interactions themselves
{[}\protect\hyperlink{ref-vasseur_environmental_2007}{39}{]}. The stochastic part of the mortality rates
(\(\epsilon_{i, t}\)) follow a normal distribution 0 centered (\(\epsilon_{i, t} \sim N(0,\sigma^2_e)\)) with a variance (\(\sigma^2_e\)). The mortality rates of
each species at the time \(t\) equals to \(d_{i,t} = d_i e^{\epsilon_{i, t}}\),
which ensured that the mortality rates never become negative
{[}\protect\hyperlink{ref-vasseur_environmental_2007}{39}{]}. We simulated \(\epsilon\) in two steps: we first
simulated a temporally correlated Ornstein-Uhlenbeck process to generate the
stochastic part of species mortality rates (\(a_{i, t}\)) and then we multiplied
those values by a correlation matrix to control the level of response
diversity. The Ornstein-Uhlenbeck process is a modified Brownian motion where
the stochastic values tend to come back to the central value (i.e.~\(a_i = 0\))
such as \(a_{i, t}\) varies according to this stochastic differential equation:
\(da_{i} = (0 - a_i)d_t + \sigma_e dW_i\), where \(dW_i\) is a Brownian motion. The
strength of environmental stochasticity was then controlled by \(\sigma_e\).

Response diversity was controlled by the correlation among species
stochastic mortality rates (\(\rho_{ij, i \neq j}\)) such that response diversity is
maximal when the stochastic component of species mortality rates (\(\epsilon_{i, t}\)) are uncorrelated (\(\rho = 0\)), and response diversity is null when the
stochastic component of species mortality rates are perfectly correlated (\(\rho = 1\)) {[}\protect\hyperlink{ref-vasseur_environmental_2007}{39},\protect\hyperlink{ref-ripa_food_2003}{41},\protect\hyperlink{ref-gouhier_synchrony_2010}{45}{]}.
Formally, the species mortality rates were a product of the stochastic mortality rates
generated by the Ornstein-Uhlenbeck process and a variance-covariance matrix
such as:

\begin{equation}
\begin{aligned}
\epsilon_{i, t} =
  \begin{bmatrix}
     \rho_{11}\sigma^2_e & \rho_{12}\sigma^2_e & \cdots & \rho_{1n}\sigma^2_e \\
     \rho_{21}\sigma^2_e & \rho_{22}\sigma^2_e & \cdots & \rho_{2n}\sigma^2_e \\
     \vdots  & \vdots  & \ddots & \vdots  \\
     \rho_{n1}\sigma^2_e & \rho_{n2}\sigma^2_e & \cdots & \rho_{nn}\sigma^2_e
  \end{bmatrix}
  \begin{bmatrix}
     a_{1,t} \\
     a_{2,t} \\
     \vdots  \\
     a_{n,t}
  \end{bmatrix}
\label{eq:rho}
\end{aligned}
\end{equation}

Where \(a_{i,t}\) is the stochastic value generated by the Ornstein-Uhlenbeck process
at time \(t\), \(\rho_{ii} = 1.0\) while all \(\rho_{ij, i \neq j}\) were equal and
control the level of response diversity. All interspecific correlations being
equal, we assume no structure in response diversity across the trophic levels,
which can contrast with the original definition that considered response
diversity as the diversity of species responses within functional groups
{[}\protect\hyperlink{ref-elmqvist_response_2003}{14}{]}.

\hypertarget{simulation-design}{%
\subsection{Simulation design}\label{simulation-design}}

We generated food web structure using the niche model {[}\protect\hyperlink{ref-williams_simple_2000}{46}{]}
with an initial species richness from \(10\) to \(60\), and
connectance from \(0.02\) to \(0.38\) (Table S1). We required
food webs without cannibalistic links and fully connected, i.e.~we discarded
food webs that contained disconnected species.

We varied the strength of environmental stochasticity by varying \(\sigma_e\) from
\(0.1\) to \(0.6\), the range of mortality rates observed in
protists {[}\protect\hyperlink{ref-vasseur_environmental_2007}{39}{]}. We decreased response diversity by
increasing \(\rho\) from \(0\) to \(1\) (i.e.~response
diversity is \(1 - \rho\)) (Fig. S1).

We varied the body mass structure of the food web to drive variation in trophic
interaction strength and mortality rates, because metabolic and mortality rates
(\(x_i\) \& \(d_i\)) scale with species body mass. To do so, we varied the
predator-prey mass ratio (\(Z\)) from \(1\) to \(100\) (Fig.
S1b). Species body masses (\(M_i\)) were computed as: \(M_i = Z^{TL_i - 1}\), where
\(TL_i\) is the trophic level of species \(i\). We set \(h = 2\), i.e.~a functional
response of type III (eq. \eqref{eq:holling}), and varied predator interference
(i.e.~from \(c = 0\) to \(c = 1\), eq. \eqref{eq:holling}),
as predator interference has been shown to have a tremendous role on population
stability {[}\protect\hyperlink{ref-brose_allometric_2006}{23}{]}.

In the absence of environmental stochasticity and with
a type III functional response, species dynamics reach a steady state without
visible oscillations. We set uniform consumer preferences for their \(n\)
resources such as: \(\omega_{ij} = 1 / n\), half-saturation rate to \(B_0=0.5\),
the maximum consumption rate to \(y_i = 4\), \(r_i=1.0\) based on previous studies
{[}\protect\hyperlink{ref-brose_allometric_2006}{23}{]}.
Intraspecific competition was set to unity (\(\alpha_{ii} = 1.0\)) and
interspecific competition (\(\alpha_{ij, i \neq j}\)) to 0.5. We set a global
carrying capacity of \(K' = 10\) to ensure that consumers were not limited by the
biomass input from primary producers. We standardised the carrying capacity by
the number of primary producers and by the interspecific competition among
producers to ensure that the effects of species richness on temporal stability
are not driven by the increase in the number of primary producers, thereby
trivially increasing biomass input for consumers. We then standardised \(K'\) such
as \(K_i = \frac{1 + \alpha_{ij} (S_{prod.} - 1)}{S_{prod.}}K'\)
{[}\protect\hyperlink{ref-ives_stability_2000}{28}{]}, where \(S_{prod.}\) is the number of primary producer
species. However, we found that our results were not affected when removing the
standardisation of carrying capacity (Fig. S2).

We numerically solved the stochastic differential equation systems during at
least 2000 timesteps and collected the last 500 timesteps, similarly to previous
studies {[}\protect\hyperlink{ref-brose_allometric_2006}{23}{]}. We numerically solve the system with \(\Delta_t = 0.1\) or \(\Delta_t = 0.05\) when instability was detected during the simulations
due to stiff changes in biomass. We set extinction threshold to
\(10^{-6}\) and if an extinction happened in the last 500 timesteps, we continued
to run the simulation for another 1000 timesteps until there was no extinction
events during the last 500 timesteps. When we found disconnected species at the
end of the simulation (i.e.~a primary producer without consumer or a consumer
without prey), we set their biomass to 0 and re-run the simulation for another
1000 timesteps. Except in the case of disconnected species, we re-ran
simulations with species starting biomass equal to their biomass at the last
timestep of the previous simulation. Supplementary analysis showed that
alternative choices of simulation processes, such as longer simulations,
rebuilding food webs (i.e.~resetting consumer preferences, recomputing species
body mass according to the new food web) after removing disconnected species did
not affect our results (Fig. S2). The food web generation and simulations were
done in Julia by developing an extension of the package
\texttt{EcologicalNetworksDynamics.jl} {[}\protect\hyperlink{ref-lajaaiti_ecologicalnetworksdynamicsjl_2025}{43}{]}.

In sum, we generated 46880
simulations over a wide range of food web structure (Median (5\%, 95\%),
species richness: 13.0 (5.0,26.0), connectance: 0.12 (0.058,0.3),
average interaction strength: 0.02 (0.0031,0.071), maximum trophic
level: 2.5 (2.0,3.5)) (Table S2).

\hypertarget{simulation-outputs}{%
\subsection{Simulation outputs}\label{simulation-outputs}}

We measured food web properties that have been linked with stability in the
literature. We measured total biomass, species richness, connectance, average
trophic level weighted by biomass, average omnivory, and average interaction
strength. Food web properties were measured at each of the 500 timesteps and
then averaged, except connectance and trophic level that were static.
Connectance was computed as \(C = L / S(S-1)\), where \(L\) is the number of trophic
links and \(S\) is the number of species in the community. Species trophic
level was computed recursively from the bottom of the food web, such as the
trophic level of a consumer was equal to the average trophic level of its
resources added to 1, the trophic level of primary producers was set to 1
{[}\protect\hyperlink{ref-christensen_ecopath_1992}{47}{]}. Average trophic level (\(\hat{wTL}\)) was then equal to:
\(\hat{wTL} = \sum_{i}1/\hat{B_i} \times \mathrm{TL}_i\), where \(\mathrm{TL}_i\) is the trophic level and
\(\hat{B_i}\) the average biomass of the species \(i\). The degree of omnivory of a
consumer was computed such as the sum of squares of its resource trophic levels
weighted by the relative preference of the consumer for each resource: \(\sum_i (\mathrm{TL}_i - \hat{\mathrm{TL}})^2 \omega_i/ \sum \omega_i\)
{[}\protect\hyperlink{ref-christensen_ecopath_1992}{47}{]}. Interaction strength was quantified as
the biomass fluxes going from every consumer/resource pair, such as: \(I_{ij} = x_i y_i B_i F_{ij}\) (eq. \eqref{eq:cons}), and averaged over time. Our way of measuring
interaction strength is very similar to previous metrics used in previous
studies with the difference here that we average fluxes over time and not the
maximal interaction strength {[}\protect\hyperlink{ref-mccann_weak_1998}{20}{]}. Finally, the average
interaction strength of a community was computed as the average of the
\(S \times S\) interaction strength matrix at the exclusion of null interaction
strengths (i.e.~excluding absence of trophic interactions).

We assessed the effects of food web structure, environmental stochasticity and
response diversity on the temporal stability of community biomass (\(S_{com}\)).
We measured temporal stability of biomass as the inverse of the coefficient of
variation. We further partitioned stability into population stability
(\(S_{pop}\)) and asynchrony (\(\phi\)) as:

\begin{equation}
\begin{aligned}
S_{com} = S_{pop} \times \phi
\label{eq:scom}
\end{aligned}
\end{equation}

with \(S_{com} = \mu_{tot} / \sigma_{tot}\), \(S_{pop} = \mu_{tot} / \sum_i \sigma_i\),
\(\phi = \sum_i \sigma_i / \sigma_{tot}\); where \(\mu_{tot}\) and \(\sigma_{tot}\)
are respectively average total biomass and standard deviation of total biomass
{[}\protect\hyperlink{ref-zhao_biodiversity_2022}{8},\protect\hyperlink{ref-thibaut_understanding_2013}{12}{]}.

To obtain further mechanistic insight about the effects of environmental
stochasticity, response diversity, and food web structure on stability, we
partitioned asynchrony into portfolio effects (\(PFE\)) and compensatory dynamics
stemming from species interactions {[}\protect\hyperlink{ref-zhao_biodiversity_2022}{8}{]} (\(CPE_{int}\)). The
measurement of portfolio effects has been a topic of debate
{[}\protect\hyperlink{ref-zhao_biodiversity_2022}{8},\protect\hyperlink{ref-loreau_biodiversity_2021}{48}{]}. Doak's definition
{[}\protect\hyperlink{ref-doak_statistical_1998}{4}{]} defines portfolio effects as the product of statistical
averaging effects (\(SAE\)) and compensatory effects arising from response
diversity {[}\protect\hyperlink{ref-doak_statistical_1998}{4},\protect\hyperlink{ref-zhao_biodiversity_2022}{8}{]} (\(CPE_{env}\)). It
then quantities the portfolio effect (\(PFE = SAE \times CPE_{env}\)) emerging
from independent species fluctuations, i.e.~the statistical averaging effect
(\(SAE\)), weighted by the effects of response diversity of species to
environmental fluctuations, such as portfolio effects are dampened if species
have correlated responses to environmental fluctuations. To compute portfolio
effects and compensatory dynamics stemming from species interactions, we i)
partition asynchrony into statistical averaging effect and compensatory dynamics
(i.e.~\(\phi = SAE \times CPE\), eq. \eqref{eq:scom2}); and ii) partition
compensatory dynamics originating from species interactions (\(CPE_{int}\)) and
from response diversity (\(CPE_{env}\)).

\begin{equation}
\begin{aligned}
S_{com} &= S_{pop} \times SAE \times CPE \\
S_{com} &= S_{pop} \times SAE \times CPE_{int} \times CPE_{env}
\label{eq:scom2}
\end{aligned}
\end{equation}

Statistical averaging effect (\(SAE\)) is the share of asynchrony assuming that
species are independent, meaning that the total variance of the community is
equal to the sum of species variances only (i.e.~species covariances are null).
We can define community stability (\(S_{com, IP}\)) in that scenario and then
derive \(SAE\):

\begin{equation}
\begin{aligned}
S_{com, IP} &= \dfrac{\mu_{tot}}{\sqrt{\sum_i \sigma^2_i}} = SAE \times S_{pop} \\
SAE &= \dfrac{S_{com, IP}}{S_{pop}} = \dfrac{\sum_i \sigma_i}{\sqrt{\sum_i \sigma^2_i}}\\
\end{aligned}
\label{eq:sae2}
\end{equation}

Compensatory dynamics in the food webs is then the remaining part of asynchrony:

\begin{equation}
\begin{aligned}
CPE &= \dfrac{S_{com}}{SAE \times S_{pop}} = \dfrac{\sqrt{\sum_i \sigma^2_i}}{\sigma_{tot}}\\
\end{aligned}
\label{eq:cpe}
\end{equation}

Compensatory effects can arise from predator-prey interactions
(\(CPE_{int}\)) or from differential response of species to environmental
stochasticity (\(CPE_{env}\)). We computed \(CPE_{env}\) on the part of
biomass fluctuations that are due to environmental stochasticity and response
diversity (\(\mathbf{B_{d, i}}\), eq. \eqref{eq:cpenv}), i.e.~from the
deterministic and stochastic parts of mortality rates. To do so, we used the
time x species matrices of species biomass \(\mathbf{B}\) and stochastic mortality
rates \(\mathbf{E}\):

\begin{equation}
\begin{aligned}
\mathbf{B_{d, i}} &= \mathbf{B} \times \mathrm{diag}(d_i) \times e^\mathbf{E}
\end{aligned}
\label{eq:cpenv}
\end{equation}

We computed \(CPE_{env}\) using the equation \eqref{eq:cpe} with the time x species
matrix (\(\mathbf{B_{d, i}}\)), so that compensatory effects due to environmental
fluctuations are mainly determined by response diversity (Fig. S3). Then
compensatory effects due to species interactions were computed as:
\(CPE_{int} = CPE / CPE_{env}\). Finally, portfolio effects (\(PFE\)) were computed
as: \(PFE = SAE \times CPE_{env}\) and compensatory effects due to species
interactions as \(CPE_{int}\). The final stability decomposition reads as follows:

\begin{equation}
\begin{aligned}
S_{com} = S_{pop} \times PFE \times CPE_{int}
\label{eq:scom3}
\end{aligned}
\end{equation}

\hypertarget{statistical-analysis}{%
\subsection{Statistical analysis}\label{statistical-analysis}}

Using a structural equation model, we estimate the contribution to temporal
community stability of the numerous processes reviewed in the introduction
across the multiple ecological scales at which they occur (Fig 1).

The arrangement of the structural equation model was defined a priori based on
the logic of what ecological scales each process operates on and on
expectations from previous theoretical studies. Working from right to left in
Fig. \ref{fig:sem}, we first specify response diversity and environmental stochasticity
which operate across all species in a community. Richness, average interaction
strength, connectance and average trophic level then represent structural
properties of the community (food web). Moving further left, the portfolio
effect and compensatory effects due to species interactions represent
aggregated species specific contributions to processes that underpin
asynchrony. Here, asynchrony is paired with population stability,
also an aggregated process at the species scale.

The flow from right to left represents expectations from theory. Theory
suggests a strong influence of average trophic level, average interaction
strength, and connectance on population stability and asynchrony
{[}\protect\hyperlink{ref-mccann_weak_1998}{20}--\protect\hyperlink{ref-fahimipour_compensation_2017}{22},\protect\hyperlink{ref-mccann_reevaluating_1997}{27}{]}, directly for population stability but mediated by portfolio and
compensatory effects for asynchrony {[}\protect\hyperlink{ref-zhao_biodiversity_2022}{8}{]}. The SEM then reflects expected
impacts of environmental stochasticity and response diversity on population
stability and asynchrony partitions. Furthest left, our SEM also captures the
core partition of asynchrony into compensatory effects generated by species
interactions (\(CPE_{int}\)) and portfolio effects stemming from statistical averaging
and compensatory effects due to response diversity (\(SAE \times CPE_{env}\)), and from
asynchrony and population stability to community stability.

Though not shown in Fig. \ref{fig:sem}, we included the predator-prey mass ratio and
predator interference values as control variables as they have been shown to
influence population stability {[}\protect\hyperlink{ref-brose_allometric_2006}{23}{]}. And while average omnivory is also
expected to increase population stability {[}\protect\hyperlink{ref-mccann_reevaluating_1997}{27}{]}, it was
not included because of strong collinearity (i.e.~measured by the Variance
Inflation Factor) with average trophic level.

Prior to evaluating hypotheses in the structural equation model, we logged all
the stability components to transform their relation from multiplicative to
additive (eq. \eqref{eq:scom2}). We ensured that all the linear models composing
the structural equation model presented low multicollinearity (Variance
Inflation Factor \textless{} 3, Table S3), although interaction strength and species
richness were highly correlated (Fig. S4). The structural equation model was
evaluated using the R package \texttt{PiecewiseSEM} {[}\protect\hyperlink{ref-lefcheck_piecewisesem_2016}{49}{]}. We
computed the sum of the direct and indirect effects using \texttt{semEff} R package
{[}\protect\hyperlink{ref-murphy_semeff_2021}{50}{]}. In the main text, we reported standardised coefficients
which were obtained by scaling the coefficients by the standard deviation of
the response and predictor variable. Finally, we reported only the direct
standardised coefficients with an absolute value equal or above \(0.05\), because
almost all effects were statistically significant given the high numbers of
simulations. These direct standardised coefficients reflect partial correlation
coefficients, while total standardised coefficients reflect correlation
coefficients {[}\protect\hyperlink{ref-grace_structural_2006}{51}{]}.

In the second analysis, we used linear mixed effect models to test how food web
structure, response diversity, and environmental stochasticity modulate
stability-species richness relationships. We modelled temporal stability of
community biomass according to food web structure: species richness, weighted
average trophic level, connectance and average interaction strength, as well as
environmental stochasticity and response diversity. We further included the
two-way interactions between species richness and food web structure, between
species richness and response diversity, between species richness and
environmental stochasticity. Finally, we added the three-way interactions
between species richness, food web structure and response diversity, between
species richness, food web structure and environmental stochasticity. We
included predator interference, predator-prey mass ratio as control variables
and added the ID of the food web as a random intercept.

We used this model to predict the shape of the relationship between community
stability and species richness along a gradient of response diversity and of
food web structure. To do so, we summed all the coefficients involving species
richness in the model, i.e.~the one, two, and three-way terms. The values of
the other variables were set to the values indicated in Fig. \ref{fig:pred}. The
coefficients of the models are displayed in Fig. S5. The VIFs were lower than
3, indicating low multicollinearity (Table S4). We checked the distribution of
the residuals (Fig. S6) using \texttt{DHARMa} R package. The model was implemented using
the R package \texttt{glmmTMB} {[}\protect\hyperlink{ref-brooks_glmmtmb_2017}{52}{]}.

\hypertarget{results}{%
\section{Results}\label{results}}

\begin{figure}
\centering
\includegraphics{p_sem_fig.png}
\caption{\label{fig:sem}Structural Equation Model linking food web structure, environmental stochasticity, and response diversity to the temporal stability of community biomass. a, Continuous blue arrows display positive and dashed red arrows display negative standardised effects, the arrow widths are proportional to the absolute values of the standardised effects. Only the standardised coefficients whose absolute values are greater to 0.05 are displayed. The full table of coefficients including control variables is displayed in Table S5. b, Total effects of food web structure, environmental stochasticity, and response diversity on community stability, asynchrony, and population stability derived from the structural equation model. N = 46880 simulations.}
\end{figure}

\hypertarget{the-drivers-of-community-stability}{%
\subsection{The drivers of community stability}\label{the-drivers-of-community-stability}}

Applying a structural equation model on the outcomes of food web simulations
revealed that the temporal stability of community biomass is positively
influenced by both population stability and asynchrony, and that population
stability has nearly double the contribution of asynchrony (resp. \(r_\delta\)=
0.9 and 0.57, \(r_\delta\) is the standardised coefficient, Fig.
\ref{fig:sem}a). This result might be driven
by the relative variance of environmental stochasticity and response diversity
explored in our simulations, since a higher variance in environmental
stochasticity relative to response diversity would result in a higher effect
size in population stability than in asynchrony because those coefficients were
expressed relative to the standard deviation of the variables (i.e.
standardised coefficients). However, the standard deviation of environmental
stochasticity was two times lower than response diversity ( resp.
0.17 vs
0.35), so that the higher importance of
population stability relative to asynchrony that we report cannot be explained
by the explored parameter ranges in the study. All the tested causal
relationships are presented in Table S5.

We further found that asynchrony is driven more by portfolio effects than by
compensatory dynamics (resp. \(r_\delta\)= 0.92 and 0.36). In more details, response
diversity has strong positive effects on asynchrony, mainly through portfolio
effects (\(r_\delta\)= 0.85), at the expense of compensatory dynamics
(\(r_\delta\)= -0.12). All other studied properties were found to have much
weaker effects on asynchrony. Food web structural properties -- namely average
trophic level, connectance and average interaction strength -- have consistent
negative effects on asynchrony,
through both portfolio and compensatory effects (Fig. \ref{fig:sem}a). Indeed, as average
trophic level and connectance increase, compensatory effects are reduced (resp.
\(r_\delta\)= -0.17 and -0.16, Fig. \ref{fig:sem}a), while average
interaction strength has a small
negative effect on portfolio effects (\(r_\delta\)= -0.05). Species richness has a
positive effect on both portfolio and compensatory effects (resp. \(r_\delta\)= 0.18 and
0.14). Finally, as expected by the standardisation of asynchrony to total
community variance (See Methods), environmental stochasticity has little effect
on asynchrony.

Focusing on population stability, environmental stochasticity was found to have
a strong negative effect on it (\(r_\delta\)= -0.91). All other studied properties were
found to have little effect on population stability, except for the average
trophic level which has a small positive effect (\(r_\delta\)= 0.17, Fig.
\ref{fig:sem}a, but see Fig. S2) and response diversity a small negative effect
(\(r_\delta\)= -0.07).

Deriving total effects from the sum of direct and indirect effects in the
structural equation model (See Methods), we showed that both environmental
stochasticity and species' response diversity have by far the strongest total
effects on community stability (resp. \(r_\delta\)= -0.8 and 0.36,
Fig. \ref{fig:sem}b). Species richness and average trophic level have similar
small but positive total effects on community stability (\(r_\delta\)= 0.08),
but through opposite pathways. While species richness increases community
stability through a strong positive effect on asynchrony (\(r_\delta\)= 0.21),
it also weakly decreases population stability (\(r_\delta\)= -0.04). In
contrast, average trophic level increases community stability through strong
positive effect on population stability (\(r_\delta\)= 0.17), but is dampened
by a negative total effect on asynchrony (\(r_\delta\)= -0.1). Interestingly,
average interaction strength also has a positive effect on population stability
and a negative one on asynchrony (resp. \(r_\delta\)= 0.02 and -0.06). This
results in a total negative effect on community stability (\(r_\delta\)=
-0.02). Connectance, instead, has a negative effect on community stability
(\(r_\delta\)= -0.06) through negative effects on both asynchrony and
population stability (resp. \(r_\delta\)= -0.07 and -0.02).

\begin{figure}
\centering
\includegraphics{p_fig3.png}
\caption{\label{fig:pred}Slope of the community stability - species richness relationship based on food web structure and response diversity. a, Diagrams displaying the slope coefficients of the effect of species richness on community stability. The average values of the control variables were used to generate the predictions: connectance = 0.14, average interaction strength = 0.025, average trophic level = 1.49, predator interference = 0.5, Predator-Prey Mass Ratio = 33, and environmental stochasticity = 0.3. b, c: Examples of predictions from the model for two combinations of average trophic level and response diversity (values displayed in panel a). Lines display the mean predictions of the model. The coefficients of the linear model are displayed in Fig. S4. Relationships between population stability, asynchrony and species richness are displayed in Fig. S7.}
\end{figure}

\hypertarget{the-context-dependence-of-stability---richness-relationships}{%
\subsection{The context dependence of stability - richness relationships}\label{the-context-dependence-of-stability---richness-relationships}}

We further assessed how environmental stochasticity and response diversity
interacted with food web structural properties to determine the sign of the
stability-richness relationship. We found that there are strong interactions
between response diversity and food web structure, which jointly determine the
sign of the stability-richness relationship.

In the absence of response diversity, we found only negative stability-richness
relationships, regardless of the food web structure (Fig. \ref{fig:pred}a). Response
diversity enables the rise of positive stability-richness relationships,
enhanced by average interaction strength and dampened by average trophic level
and connectance. Higher interaction strength leads to positive
stability-richness relationships for lower values of response diversity, and
interaction strength enhances asynchrony-richness relationships (Fig. S7), so
that the stronger positive stability-richness relationships are found at both
high interaction strength and high response diversity levels. In contrast,
higher average trophic level and connectance both lead to negative
stability-richness relationships, because they drive more negative population
stability-richness relationships and weaker positive asynchrony-richness
relationships (Fig. S7).

Overall, our results highlight a contrast between the relatively small effects
of food web structural properties on overall community stability (Fig.
\ref{fig:sem}b) but their strong effects on the slope of the
stability-richness relationship (Fig. \ref{fig:pred}a). A contrast which is
consistent when considering interactive effects among food web structure,
response diversity and environmental stochasticity (Fig. S3).

\hypertarget{discussion}{%
\section{Discussion}\label{discussion}}

With projected changes in climate and land use, and associated loss of
diversity, change in community structure, potential homogenisation, it is
important to understand the mechanisms which allow ecological communities to
remain stable despite perturbations. Recent research has focused on either the
average stability of the populations composing the community or on the
asynchrony of the population fluctuations {[}\protect\hyperlink{ref-zhao_biodiversity_2022}{8},\protect\hyperlink{ref-thibaut_understanding_2013}{12}{]}. Moreover, the majority of
this work has been done in simplified assemblages, such as single trophic level
communities. In this focused and often simplified context community stability
has been shown to be driven by asynchrony in species fluctuations, rather than
by population stability {[}\protect\hyperlink{ref-olivier_urbanization_2020}{13}{]}.

However, whether this inference holds for more complex, trophically structured
ecological communities such as food webs, remains largely unknown. Here, we
tackled this question using simulations of complex, stochastic food web
dynamics via a bioenergetic model. In doing so, we reveal not only the relative
contribution to temporal stability of multiple processes spanning multiple
ecological scales, but also predictions about the context dependency of
diversity - stability relationships amidst rapid global environmental change.

Our results show that, in complex food webs, community stability is more driven
by population stability than by asynchrony, in contrast with what was observed
in single trophic level communities. As we showed above, this was not due to
the ratio of variance between environmental stochasticity and response
diversity. Our study further reveals that the main drivers of community
stability are environmental stochasticity and response diversity, respectively
acting on population stability and asynchrony. Furthermore, despite previous
work in simple communities suggesting that food web structural properties drive
stability, we found that these properties have much weaker effects in more
complex food webs, a result consistent across a diversity of methodological
choices (Fig. S2).

However, structural properties are central to our understanding of
stability-richness relationships. They were found to play an important role in
the sign of the stability-richness relationship, which is strongly mediated by
the interaction between food web structure and response diversity. Our results
highlight the dual importance of response diversity to improve our
understanding of stability per se and of richness-stability relationships in
structured ecological communities facing variable environments. In the
following sections, we deliver more nuanced insights into these conclusions.

\hypertarget{environmental-stochasticity-and-response-diversity-are-the-main-drivers-of-community-stability}{%
\subsection{Environmental Stochasticity and response diversity are the main drivers of community stability}\label{environmental-stochasticity-and-response-diversity-are-the-main-drivers-of-community-stability}}

We found environmental stochasticity and response diversity to be the main
overall drivers of community stability. While environmental stochasticity
decreases population stability, response diversity generates asynchrony,
thereby buffering environmental stochasticity. This result is in accordance
with previous theoretical and empirical findings in competitive communities
{[}\protect\hyperlink{ref-de_mazancourt_predicting_2013}{6},\protect\hyperlink{ref-thibaut_diversity_2012}{7},\protect\hyperlink{ref-thibaut_understanding_2013}{12},\protect\hyperlink{ref-ives_stability_2000}{28},\protect\hyperlink{ref-thebault_trophic_2005}{29},\protect\hyperlink{ref-loreau_species_2008}{34}{]}. In
turn, response diversity was found to operate on asynchrony. This is in
agreement with empirical findings in competitive assemblages where response
diversity was quantified with trait-based or non trait-based approaches
{[}\protect\hyperlink{ref-thibaut_diversity_2012}{7},\protect\hyperlink{ref-zhao_biodiversity_2022}{8},\protect\hyperlink{ref-sasaki_species_2019}{15}{]}.

Specifically, in our model, response diversity was driving portfolio effects
through compensatory effects due to species' differential responses to
environmental fluctuations (i.e.~\(CPE_{env}\), see Methods eq.
\eqref{eq:cpe}-\eqref{eq:scom3}), thereby modulating the statistical averaging of
independent species fluctuations {[}\protect\hyperlink{ref-doak_statistical_1998}{4}{]}. This mechanism
highlights how differences in species' environmental responses contribute to
asynchrony, stabilizing communities by reducing the synchrony of population
fluctuations.

A consequence of the large effect of response diversity is that asynchrony in
species fluctuations is driven more by the statistical averaging of independent
species fluctuations weighted by response diversity (portfolio effects) than by
compensatory dynamics emerging from species interactions. This finding is also
coherent with previous studies in competitive communities
{[}\protect\hyperlink{ref-thibaut_diversity_2012}{7},\protect\hyperlink{ref-zhao_biodiversity_2022}{8}{]}, but it was less expected to
happen in food webs. Trophic interactions are indeed expected to result in more
compensatory dynamics resulting from oscillations between predators and prey
{[}\protect\hyperlink{ref-vandermeer_coupled_2004}{53}{]} or prey switching from predators {[}\protect\hyperlink{ref-mccann_weak_1998}{20},\protect\hyperlink{ref-mccann_diversitystability_2000}{21}{]}. However, pioneering theoretical studies have
shown that environmental stochasticity can dampen prey switching
{[}\protect\hyperlink{ref-vasseur_environmental_2007}{39}{]} and that trophic cascades between trophic levels
in species rich food webs are weak compared to food chains and species poor
food web modules {[}\protect\hyperlink{ref-fahimipour_compensation_2017}{22}{]}. So, the question of the
importance of compensatory dynamics in complex food webs remains an open and
interesting avenue for future research.

Population stability (not asynchrony) was found to be the dominant driver of
community stability in our food web simulations, through the large effect of
environmental stochasticity. These results are in agreement with previous
findings in empirical food webs {[}\protect\hyperlink{ref-danet_species_2021}{25}{]} and food web models
{[}\protect\hyperlink{ref-eschenbrenner_diversity_2023}{26}{]}, but contrast with empirical evidence in
competitive communities {[}\protect\hyperlink{ref-olivier_urbanization_2020}{13}{]}. It suggests fundamental
differences between competitive communities and food webs in the relative
contributions of population stability and asynchrony to community stability.
Our findings suggest that the presence of trophic links contributes to
synchronise the species of the community because we found that the food web
metrics studied -- namely average trophic level, connectance, and average
interaction strength -- all decrease asynchrony. These results are in line with
empirical evidence that predators can synchronise their prey
{[}\protect\hyperlink{ref-raimondo_interspecific_2004}{18},\protect\hyperlink{ref-korpimaki_predatorinduced_2005}{19}{]}, and that
species pairs involved in trophic interactions are more synchronous than those
that are not {[}\protect\hyperlink{ref-danet_species_2021}{25}{]}.

A complementary explanation for the relatively lower importance of asynchrony
in food webs (compared to single trophic level communities) could be related to
the biomass structure of food webs. Food webs typically include a larger range
of body masses than single trophic level assemblages {[}\protect\hyperlink{ref-brose_predator_2019}{54}{]}.
Such highly uneven biomass distribution can translate in uneven species
temporal variability, thereby creating uneven species contributions to
asynchrony {[}\protect\hyperlink{ref-thibaut_understanding_2013}{12}{]} and limiting portfolio effects and
compensatory dynamics in species-rich communities {[}\protect\hyperlink{ref-zhao_biodiversity_2022}{8}{]}.

Investigating in more detail the drivers of community stability, our results
suggest that species richness and the network structural properties
investigated only have a weak effect on overall community stability. We found
that a higher average trophic level leads to higher community stability by
increasing the average population stability, confirming previous theoretical
{[}\protect\hyperlink{ref-brose_allometric_2006}{23},\protect\hyperlink{ref-shanafelt_stability_2018}{24},\protect\hyperlink{ref-eschenbrenner_diversity_2023}{26}{]} and empirical {[}\protect\hyperlink{ref-danet_species_2021}{25}{]} results. However,
this effect disappeared when we simulated community dynamics with equal
mortality rates for all species instead of allometric ones (Fig. S2). This
suggests that the higher stability of food webs with higher average trophic
levels stems more from their lower mortality rates rather than from their
trophic role directly. The overall effect of connectance on overall community
stability was found to be quite low, in agreement with previous empirical and
theoretical findings {[}\protect\hyperlink{ref-danet_species_2021}{25},\protect\hyperlink{ref-eschenbrenner_diversity_2023}{26}{]}.
Similarly, we found very low effects of the average interaction strength on
community stability, which is in apparent disagreement with the predominant
role of interaction strength and connectance on stability {[}\protect\hyperlink{ref-may_will_1972}{2}{]}. We
found that average interaction strength had a weak positive effect on
population stability, in confirmation with previous theoretical findings
aligning with our model assumption that there was no asymmetry of consumer
preferences for their prey {[}\protect\hyperlink{ref-gellner_consistent_2016}{55}{]}.

\hypertarget{stability-richness-relationships-depend-on-response-diversity-and-food-web-structure}{%
\subsection{Stability-richness relationships depend on response diversity and food web structure}\label{stability-richness-relationships-depend-on-response-diversity-and-food-web-structure}}

Our results also allow us to revisit the old question of the relationship
between the complexity of ecological communities and their stability
{[}\protect\hyperlink{ref-may_will_1972}{2},\protect\hyperlink{ref-angelis_stability_1975}{56}{]}. We found that response diversity
enables the emergence of positive stability-richness relationships, enhanced by
average interaction strength and dampened by average trophic level and
connectance. In the absence of response diversity, only negative relationships
between community stability and species richness are observed. This reinforces
the idea that increases in species richness enhance community stability only if
the added species have different responses to environmental perturbations, i.e.
if there is response diversity {[}\protect\hyperlink{ref-de_mazancourt_predicting_2013}{6},\protect\hyperlink{ref-thebault_trophic_2005}{29}{]}. Indeed, a community with numerous species which
respond differently to environmental change is more likely to buffer a wide
array of perturbations {[}\protect\hyperlink{ref-yachi_biodiversity_1999}{3},\protect\hyperlink{ref-ives_stability_2000}{28}{]}.

Despite its relatively weaker effect on community stability (than response
diversity and environmental stochasticity), food web structure interacts
strongly with response diversity to determine the sign of stability-richness
relationships. Surprisingly, we find that higher average interaction strength
enhances positive stability-richness relationships by enhancing
asynchrony-richness relationships. Conversely, higher average trophic level and
higher connectance both led to negative stability-richness relationships by
dampening asynchrony-richness relationships. It would appear that asynchrony
becomes too weak to compensate for the simultaneous decrease in population
stability induced by increasing species richness.

These results should be understood as the predicted stability-richness
relationships if food web structure is kept constant, i.e.~if food web
structure does not vary with species richness. However, connectance
{[}\protect\hyperlink{ref-dunne_network_2006}{36}{]} and average trophic level weighted by biomass are
documented to decrease with species richness {[}\protect\hyperlink{ref-danet_species_2021}{25}{]}, because of
energetic and topological constraints {[}\protect\hyperlink{ref-dunne_network_2006}{36},\protect\hyperlink{ref-gravel_trophic_2011}{57}{]}. Our results add to this body of work by suggesting that
communities increasing in species richness without concurrent decreases in
connectance and average trophic level should experience reduced stability. Our
results resonate with findings that some community structure can lead to
negative stability-richness relationships, as previously reported in freshwater
food webs {[}\protect\hyperlink{ref-danet_species_2021}{25}{]}. Response diversity might therefore be one of
the mechanisms that explains why positive relationships between stability and
species richness are so prevalent in empirical settings
{[}\protect\hyperlink{ref-thibaut_understanding_2013}{12},\protect\hyperlink{ref-elmqvist_response_2003}{14}{]}, while food web
theoretical models that do not often include environmental stochasticity and
response diversity typically find negative complexity-stability relationships
{[}\protect\hyperlink{ref-ives_stability_2007}{40}{]}.

\hypertarget{summary}{%
\subsection{Summary}\label{summary}}

In conclusion, our study highlights the contrasting effects of response
diversity and food web structure on the temporal stability of species rich
communities: response diversity underpins community stability and strongly
interacts with food web structure to define stability-richness relationships.
In line with the insurance hypothesis {[}\protect\hyperlink{ref-yachi_biodiversity_1999}{3},\protect\hyperlink{ref-loreau_biodiversity_2021}{48}{]}, our results suggest that environmental
stochasticity and response diversity are the main mechanisms driving community
stability, while food web structural properties are only mediating them. This
further aligns with the idea that, in species-rich communities, community
structure might be secondary for understanding the generic effects of
biodiversity on ecosystem functioning {[}\protect\hyperlink{ref-ives_stability_2000}{28},\protect\hyperlink{ref-barbier_generic_2018}{58}{]}. Our results further add to the mounting evidence that
population stability is a stronger driver of community stability than
asynchrony in food webs, in contrast with what was observed in competitive
communities.

The fact that environmental stochasticity acting with response diversity was
found to be the main driver of overall community stability raises concerns
about the consequences of the current rise in the frequency and severity of
environmental extreme events for ecological communities
{[}\protect\hyperlink{ref-ghosh_temperature_2024}{59},\protect\hyperlink{ref-lee_ipcc_2023}{60}{]}, which constitutes an increase in
environmental stochasticity. Indeed, our results suggest that these changes are
unlikely to be compensated by biotic mechanisms alone. Documenting how rapid
biodiversity changes {[}\protect\hyperlink{ref-dornelas_assemblage_2014}{61},\protect\hyperlink{ref-blowes_geography_2019}{62}{]} driven
by human pressures {[}\protect\hyperlink{ref-danet_past_2024}{63}{]} -- and in particular the reported
homogenisation {[}\protect\hyperlink{ref-olden_homogocene_2018}{64},\protect\hyperlink{ref-wang_biotic_2021}{65}{]} in community
composition -- affect response diversity will be key to accurately predict how
it will affect, in turn, the stability of communities.

Looking forward, our study brings evidence that asynchrony is predominantly
driven by portfolio effects as previously reported in plant communities.
However, we expect that some ecological contexts could result in the dominance
of compensatory effects due to species interactions over portfolio effects. An
example is the role of the structure of response diversity, since we here
considered species responses correlation across whole food web communities. We
expect that considering species differential responses to environmental
stochasticity within functional groups {[}\protect\hyperlink{ref-elmqvist_response_2003}{14}{]}, such as
trophic levels, may give more importance to compensatory dynamics due to
species interactions relative to the ones related to portfolio effects.
The structure of response diversity in food webs has been overlooked so far,
but this knowledge is expected to rise with the methodological progress in
measuring response diversity {[}\protect\hyperlink{ref-ross_how_2023}{16},\protect\hyperlink{ref-polazzo_measuring_2024}{17}{]}. Another
context to explore is the effect of the distribution of interaction strengths
in food webs. Implementing asymmetric distributions of consumer preferences for
their prey could also increase the importance of compensatory dynamics due to
species interactions as strong asymmetry in interaction strength can enhance
prey switching behaviour {[}\protect\hyperlink{ref-mccann_weak_1998}{20},\protect\hyperlink{ref-mccann_diversitystability_2000}{21}{]}.

\hypertarget{acknowledgements}{%
\section{Acknowledgements}\label{acknowledgements}}

We thank Owen Petchey, Tanya Strydom, and the anonymous reviewers for their
constructive feedback on the study. AD, TFJ and APB have received support from
the Natural Environment Research Council under grant agreement no. NE/T003502/1
awarded to APB.

\hypertarget{data-accessibility-statement}{%
\section{Data accessibility statement}\label{data-accessibility-statement}}

The code used to run the simulations, analyse them and to write the manuscript
is archived on Zenodo: \url{https://zenodo.org/records/16745483}.

\hypertarget{statement-of-authorship}{%
\section{Statement of authorship}\label{statement-of-authorship}}

AD and APB designed the research, AD performed the research. AD wrote the first
draft of the manuscript. All authors contributed substantially to revisions.

\hypertarget{references}{%
\section*{References}\label{references}}
\addcontentsline{toc}{section}{References}

\hypertarget{refs}{}
\begin{CSLReferences}{0}{0}
\leavevmode\vadjust pre{\hypertarget{ref-macarthur_fluctuations_1955}{}}%
\CSLLeftMargin{1. }%
\CSLRightInline{MacArthur R. 1955 Fluctuations of {Animal} {Populations} and a {Measure} of {Community} {Stability}. \emph{Ecology} \textbf{36}, 533--536. (doi:\href{https://doi.org/10.2307/1929601}{10.2307/1929601})}

\leavevmode\vadjust pre{\hypertarget{ref-may_will_1972}{}}%
\CSLLeftMargin{2. }%
\CSLRightInline{May RM. 1972 Will a {Large} {Complex} {System} be {Stable}? \emph{Nature} \textbf{238}, 413--414. (doi:\href{https://doi.org/10.1038/238413a0}{10.1038/238413a0})}

\leavevmode\vadjust pre{\hypertarget{ref-yachi_biodiversity_1999}{}}%
\CSLLeftMargin{3. }%
\CSLRightInline{Yachi S, Loreau M. 1999 Biodiversity and ecosystem productivity in a fluctuating environment: {The} insurance hypothesis. \emph{Proceedings of the National Academy of Sciences} \textbf{96}, 1463--1468. (doi:\href{https://doi.org/10.1073/pnas.96.4.1463}{10.1073/pnas.96.4.1463})}

\leavevmode\vadjust pre{\hypertarget{ref-doak_statistical_1998}{}}%
\CSLLeftMargin{4. }%
\CSLRightInline{Doak DF, Bigger D, Harding EK, Marvier MA, O'Malley RE, Thomson D. 1998 The {Statistical} {Inevitability} of {Stability}‐{Diversity} {Relationships} in {Community} {Ecology}. \emph{The American Naturalist} \textbf{151}, 264--276. (doi:\href{https://doi.org/10.1086/286117}{10.1086/286117})}

\leavevmode\vadjust pre{\hypertarget{ref-tilman_ecological_1999}{}}%
\CSLLeftMargin{5. }%
\CSLRightInline{Tilman D. 1999 The ecological consequences of changes in biodiversity: A search for general principles. \emph{Ecology} \textbf{80}, 1455--1474. (doi:\href{https://doi.org/10.1890/0012-9658(1999)080\%5B1455:TECOCI\%5D2.0.CO;2}{10.1890/0012-9658(1999)080{[}1455:TECOCI{]}2.0.CO;2})}

\leavevmode\vadjust pre{\hypertarget{ref-de_mazancourt_predicting_2013}{}}%
\CSLLeftMargin{6. }%
\CSLRightInline{De Mazancourt C \emph{et al.} 2013 Predicting ecosystem stability from community composition and biodiversity. \emph{Ecology Letters} \textbf{16}, 617--625. (doi:\href{https://doi.org/10.1111/ele.12088}{10.1111/ele.12088})}

\leavevmode\vadjust pre{\hypertarget{ref-thibaut_diversity_2012}{}}%
\CSLLeftMargin{7. }%
\CSLRightInline{Thibaut LM, Connolly SR, Sweatman HPA. 2012 Diversity and stability of herbivorous fishes on coral reefs. \emph{Ecology} \textbf{93}, 891--901. (doi:\href{https://doi.org/10.1890/11-1753.1}{10.1890/11-1753.1})}

\leavevmode\vadjust pre{\hypertarget{ref-zhao_biodiversity_2022}{}}%
\CSLLeftMargin{8. }%
\CSLRightInline{Zhao L, Wang S, Shen R, Gong Y, Wang C, Hong P, Reuman DC. 2022 Biodiversity stabilizes plant communities through statistical-averaging effects rather than compensatory dynamics. \emph{Nature Communications} \textbf{13}, 7804. (doi:\href{https://doi.org/10.1038/s41467-022-35514-9}{10.1038/s41467-022-35514-9})}

\leavevmode\vadjust pre{\hypertarget{ref-downing_ecosystem_2002}{}}%
\CSLLeftMargin{9. }%
\CSLRightInline{Downing AL, Leibold MA. 2002 Ecosystem consequences of species richness and composition in pond food webs. \emph{Nature} \textbf{416}, 837--841. (doi:\href{https://doi.org/10.1038/416837a}{10.1038/416837a})}

\leavevmode\vadjust pre{\hypertarget{ref-srednick_understanding_2024}{}}%
\CSLLeftMargin{10. }%
\CSLRightInline{Srednick G, Swearer SE. 2024 Understanding diversity--synchrony--stability relationships in multitrophic communities. \emph{Nature Ecology \& Evolution} \textbf{8}, 1259--1269. (doi:\href{https://doi.org/10.1038/s41559-024-02419-3}{10.1038/s41559-024-02419-3})}

\leavevmode\vadjust pre{\hypertarget{ref-thebault_stability_2010}{}}%
\CSLLeftMargin{11. }%
\CSLRightInline{Thébault E, Fontaine C. 2010 Stability of {Ecological} {Communities} and the {Architecture} of {Mutualistic} and {Trophic} {Networks}. \emph{Science} \textbf{329}, 853--856. (doi:\href{https://doi.org/10.1126/science.1188321}{10.1126/science.1188321})}

\leavevmode\vadjust pre{\hypertarget{ref-thibaut_understanding_2013}{}}%
\CSLLeftMargin{12. }%
\CSLRightInline{Thibaut LM, Connolly SR. 2013 Understanding diversity--stability relationships: Towards a unified model of portfolio effects. \emph{Ecology Letters} \textbf{16}, 140--150. (doi:\href{https://doi.org/10.1111/ele.12019}{10.1111/ele.12019})}

\leavevmode\vadjust pre{\hypertarget{ref-olivier_urbanization_2020}{}}%
\CSLLeftMargin{13. }%
\CSLRightInline{Olivier T, Thébault E, Elias M, Fontaine B, Fontaine C. 2020 Urbanization and agricultural intensification destabilize animal communities differently than diversity loss. \emph{Nature Communications} \textbf{11}, 2686. (doi:\href{https://doi.org/10.1038/s41467-020-16240-6}{10.1038/s41467-020-16240-6})}

\leavevmode\vadjust pre{\hypertarget{ref-elmqvist_response_2003}{}}%
\CSLLeftMargin{14. }%
\CSLRightInline{Elmqvist T, Folke C, Nyström M, Peterson G, Bengtsson J, Walker B, Norberg J. 2003 Response diversity, ecosystem change, and resilience. \emph{Frontiers in Ecology and the Environment} \textbf{1}, 488--494. (doi:\href{https://doi.org/10.1890/1540-9295(2003)001\%5B0488:RDECAR\%5D2.0.CO;2}{10.1890/1540-9295(2003)001{[}0488:RDECAR{]}2.0.CO;2})}

\leavevmode\vadjust pre{\hypertarget{ref-sasaki_species_2019}{}}%
\CSLLeftMargin{15. }%
\CSLRightInline{Sasaki T, Lu X, Hirota M, Bai Y. 2019 Species asynchrony and response diversity determine multifunctional stability of natural grasslands. \emph{Journal of Ecology} \textbf{107}, 1862--1875. (doi:\href{https://doi.org/10.1111/1365-2745.13151}{10.1111/1365-2745.13151})}

\leavevmode\vadjust pre{\hypertarget{ref-ross_how_2023}{}}%
\CSLLeftMargin{16. }%
\CSLRightInline{Ross SRP-J, Petchey OL, Sasaki T, Armitage DW. 2023 How to measure response diversity. \emph{Methods in Ecology and Evolution} \textbf{14}, 1150--1167. (doi:\href{https://doi.org/10.1111/2041-210X.14087}{10.1111/2041-210X.14087})}

\leavevmode\vadjust pre{\hypertarget{ref-polazzo_measuring_2024}{}}%
\CSLLeftMargin{17. }%
\CSLRightInline{Polazzo F, Limberger R, Pennekamp F, Ross SRP-J, Simpson GL, Petchey OL. 2024 Measuring the {Response} {Diversity} of {Ecological} {Communities} {Experiencing} {Multifarious} {Environmental} {Change}. \emph{Global Change Biology} \textbf{30}, e17594. (doi:\href{https://doi.org/10.1111/gcb.17594}{10.1111/gcb.17594})}

\leavevmode\vadjust pre{\hypertarget{ref-raimondo_interspecific_2004}{}}%
\CSLLeftMargin{18. }%
\CSLRightInline{Raimondo S, Turcáni M, Patoèka J, Liebhold AM. 2004 Interspecific synchrony among foliage-feeding forest {Lepidoptera} species and the potential role of generalist predators as synchronizing agents. \emph{Oikos} \textbf{107}, 462--470. (doi:\href{https://doi.org/10.1111/j.0030-1299.2004.13449.x}{10.1111/j.0030-1299.2004.13449.x})}

\leavevmode\vadjust pre{\hypertarget{ref-korpimaki_predatorinduced_2005}{}}%
\CSLLeftMargin{19. }%
\CSLRightInline{Korpimäki E, Norrdahl K, Huitu O, Klemola T. 2005 Predator--induced synchrony in population oscillations of coexisting small mammal species. \emph{Proceedings of the Royal Society B: Biological Sciences} \textbf{272}, 193--202. (doi:\href{https://doi.org/10.1098/rspb.2004.2860}{10.1098/rspb.2004.2860})}

\leavevmode\vadjust pre{\hypertarget{ref-mccann_weak_1998}{}}%
\CSLLeftMargin{20. }%
\CSLRightInline{McCann K, Hastings A, Huxel GR. 1998 Weak trophic interactions and the balance of nature. \emph{Nature} \textbf{395}, 794--798. (doi:\href{https://doi.org/10.1038/27427}{10.1038/27427})}

\leavevmode\vadjust pre{\hypertarget{ref-mccann_diversitystability_2000}{}}%
\CSLLeftMargin{21. }%
\CSLRightInline{McCann KS. 2000 \href{https://www.nature.com/articles/35012234}{The diversity--stability debate}. \emph{Nature} \textbf{405}, 228--233.}

\leavevmode\vadjust pre{\hypertarget{ref-fahimipour_compensation_2017}{}}%
\CSLLeftMargin{22. }%
\CSLRightInline{Fahimipour AK, Anderson KE, Williams RJ. 2017 Compensation masks trophic cascades in complex food webs. \emph{Theoretical Ecology} \textbf{10}, 245--253. (doi:\href{https://doi.org/10.1007/s12080-016-0326-8}{10.1007/s12080-016-0326-8})}

\leavevmode\vadjust pre{\hypertarget{ref-brose_allometric_2006}{}}%
\CSLLeftMargin{23. }%
\CSLRightInline{Brose U, Williams RJ, Martinez ND. 2006 Allometric scaling enhances stability in complex food webs. \emph{Ecology Letters} \textbf{9}, 1228--1236. (doi:\href{https://doi.org/10.1111/j.1461-0248.2006.00978.x}{10.1111/j.1461-0248.2006.00978.x})}

\leavevmode\vadjust pre{\hypertarget{ref-shanafelt_stability_2018}{}}%
\CSLLeftMargin{24. }%
\CSLRightInline{Shanafelt DW, Loreau M. 2018 Stability trophic cascades in food chains. \emph{Royal Society Open Science} \textbf{5}, 180995. (doi:\href{https://doi.org/10.1098/rsos.180995}{10.1098/rsos.180995})}

\leavevmode\vadjust pre{\hypertarget{ref-danet_species_2021}{}}%
\CSLLeftMargin{25. }%
\CSLRightInline{Danet A, Mouchet M, Bonnaffé W, Thébault E, Fontaine C. 2021 Species richness and food-web structure jointly drive community biomass and its temporal stability in fish communities. \emph{Ecology Letters} \textbf{24}, 2364--2377. (doi:\href{https://doi.org/10.1111/ele.13857}{10.1111/ele.13857})}

\leavevmode\vadjust pre{\hypertarget{ref-eschenbrenner_diversity_2023}{}}%
\CSLLeftMargin{26. }%
\CSLRightInline{Eschenbrenner J, Thébault É. 2023 Diversity, food web structure and the temporal stability of total plant and animal biomasses. \emph{Oikos} \textbf{2023}, e08769. (doi:\href{https://doi.org/10.1111/oik.08769}{10.1111/oik.08769})}

\leavevmode\vadjust pre{\hypertarget{ref-mccann_reevaluating_1997}{}}%
\CSLLeftMargin{27. }%
\CSLRightInline{McCann K, Hastings A. 1997 Re--evaluating the omnivory--stability relationship in food webs. \emph{Proceedings of the Royal Society of London. Series B: Biological Sciences} \textbf{264}, 1249--1254. (doi:\href{https://doi.org/10.1098/rspb.1997.0172}{10.1098/rspb.1997.0172})}

\leavevmode\vadjust pre{\hypertarget{ref-ives_stability_2000}{}}%
\CSLLeftMargin{28. }%
\CSLRightInline{Ives Ar, Klug Jl, Gross K. 2000 Stability and species richness in complex communities. \emph{Ecology Letters} \textbf{3}, 399--411. (doi:\href{https://doi.org/10.1046/j.1461-0248.2000.00144.x}{10.1046/j.1461-0248.2000.00144.x})}

\leavevmode\vadjust pre{\hypertarget{ref-thebault_trophic_2005}{}}%
\CSLLeftMargin{29. }%
\CSLRightInline{Thébault E, Loreau M. 2005 Trophic {Interactions} and the {Relationship} between {Species} {Diversity} and {Ecosystem} {Stability}. \emph{The American Naturalist} \textbf{166}, E95--E114. (doi:\href{https://doi.org/10.1086/444403}{10.1086/444403})}

\leavevmode\vadjust pre{\hypertarget{ref-tilman_biodiversity_2006}{}}%
\CSLLeftMargin{30. }%
\CSLRightInline{Tilman D, Reich PB, Knops JMH. 2006 Biodiversity and ecosystem stability in a decade-long grassland experiment. \emph{Nature} \textbf{441}, 629--632. (doi:\href{https://doi.org/10.1038/nature04742}{10.1038/nature04742})}

\leavevmode\vadjust pre{\hypertarget{ref-loreau_biodiversity_2013}{}}%
\CSLLeftMargin{31. }%
\CSLRightInline{Loreau M, Mazancourt C de. 2013 Biodiversity and ecosystem stability: A synthesis of underlying mechanisms. \emph{Ecology Letters} \textbf{16}, 106--115. (doi:\href{https://doi.org/10.1111/ele.12073}{10.1111/ele.12073})}

\leavevmode\vadjust pre{\hypertarget{ref-houlahan_negative_2018}{}}%
\CSLLeftMargin{32. }%
\CSLRightInline{Houlahan JE \emph{et al.} 2018 Negative relationships between species richness and temporal variability are common but weak in natural systems. \emph{Ecology} \textbf{99}, 2592--2604. (doi:\url{https://doi.org/10.1002/ecy.2514})}

\leavevmode\vadjust pre{\hypertarget{ref-hatton_linking_2019}{}}%
\CSLLeftMargin{33. }%
\CSLRightInline{Hatton IA, Dobson AP, Storch D, Galbraith ED, Loreau M. 2019 Linking scaling laws across eukaryotes. \emph{Proceedings of the National Academy of Sciences} \textbf{116}, 21616--21622. (doi:\href{https://doi.org/10.1073/pnas.1900492116}{10.1073/pnas.1900492116})}

\leavevmode\vadjust pre{\hypertarget{ref-loreau_species_2008}{}}%
\CSLLeftMargin{34. }%
\CSLRightInline{Loreau M, Mazancourt C de. 2008 Species {Synchrony} and {Its} {Drivers}: {Neutral} and {Nonneutral} {Community} {Dynamics} in {Fluctuating} {Environments}. \emph{The American Naturalist} \textbf{172}, E48--E66. (doi:\href{https://doi.org/10.1086/589746}{10.1086/589746})}

\leavevmode\vadjust pre{\hypertarget{ref-frund_response_2013}{}}%
\CSLLeftMargin{35. }%
\CSLRightInline{Fründ J, Zieger SL, Tscharntke T. 2013 Response diversity of wild bees to overwintering temperatures. \emph{Oecologia} \textbf{173}, 1639--1648. (doi:\href{https://doi.org/10.1007/s00442-013-2729-1}{10.1007/s00442-013-2729-1})}

\leavevmode\vadjust pre{\hypertarget{ref-dunne_network_2006}{}}%
\CSLLeftMargin{36. }%
\CSLRightInline{Dunne JA. 2006 The network structure of food webs. \emph{Ecological networks: linking structure to dynamics in food webs}, 27--86.}

\leavevmode\vadjust pre{\hypertarget{ref-galiana_spatial_2018}{}}%
\CSLLeftMargin{37. }%
\CSLRightInline{Galiana N, Lurgi M, Claramunt-López B, Fortin M-J, Leroux S, Cazelles K, Gravel D, Montoya JM. 2018 The spatial scaling of species interaction networks. \emph{Nature Ecology \& Evolution} \textbf{2}, 782--790. (doi:\href{https://doi.org/10.1038/s41559-018-0517-3}{10.1038/s41559-018-0517-3})}

\leavevmode\vadjust pre{\hypertarget{ref-yodzis_body_1992}{}}%
\CSLLeftMargin{38. }%
\CSLRightInline{Yodzis P, Innes S. 1992 Body {Size} and {Consumer}-{Resource} {Dynamics}. \emph{The American Naturalist} \textbf{139}, 1151--1175. (doi:\href{https://doi.org/10.1086/285380}{10.1086/285380})}

\leavevmode\vadjust pre{\hypertarget{ref-vasseur_environmental_2007}{}}%
\CSLLeftMargin{39. }%
\CSLRightInline{Vasseur DA, Fox JW. 2007 Environmental fluctuations can stabilize food web dynamics by increasing synchrony. \emph{Ecology Letters} \textbf{10}, 1066--1074. (doi:\href{https://doi.org/10.1111/j.1461-0248.2007.01099.x}{10.1111/j.1461-0248.2007.01099.x})}

\leavevmode\vadjust pre{\hypertarget{ref-ives_stability_2007}{}}%
\CSLLeftMargin{40. }%
\CSLRightInline{Ives AR, Carpenter SR. 2007 Stability and {Diversity} of {Ecosystems}. \emph{Science} \textbf{317}, 58--62. (doi:\href{https://doi.org/10.1126/science.1133258}{10.1126/science.1133258})}

\leavevmode\vadjust pre{\hypertarget{ref-ripa_food_2003}{}}%
\CSLLeftMargin{41. }%
\CSLRightInline{Ripa J, Ives AR. 2003 Food web dynamics in correlated and autocorrelated environments. \emph{Theoretical Population Biology} \textbf{64}, 369--384. (doi:\href{https://doi.org/10.1016/S0040-5809(03)00089-3}{10.1016/S0040-5809(03)00089-3})}

\leavevmode\vadjust pre{\hypertarget{ref-rooney_homage_2006}{}}%
\CSLLeftMargin{42. }%
\CSLRightInline{Williams RJ, Brose U, Martinez ND. 2006 Homage to {Yodzis} and {Innes} 1992: {Scaling} up {Feeding}-{Based} {Population} {Dynamics} to {Complex} {Ecological} {Networks}. In \emph{From {Energetics} to {Ecosystems}: {The} {Dynamics} and {Structure} of {Ecological} {Systems}} (eds N Rooney, KS McCann, DLG Noakes), pp. 37--51. Springer Netherlands. (doi:\href{https://doi.org/10.1007/978-1-4020-5337-5_2}{10.1007/978-1-4020-5337-5\_2})}

\leavevmode\vadjust pre{\hypertarget{ref-lajaaiti_ecologicalnetworksdynamicsjl_2025}{}}%
\CSLLeftMargin{43. }%
\CSLRightInline{Lajaaiti I, Bonnici I, Kéfi S, Mayall H, Danet A, Beckerman AP, Malpas T, Delmas E. 2025 {EcologicalNetworksDynamics}.jl: {A} {Julia} package to simulate the temporal dynamics of complex ecological networks. \emph{Methods in Ecology and Evolution} \textbf{16}, 520--529. (doi:\href{https://doi.org/10.1111/2041-210X.14497}{10.1111/2041-210X.14497})}

\leavevmode\vadjust pre{\hypertarget{ref-brown_toward_2004}{}}%
\CSLLeftMargin{44. }%
\CSLRightInline{Brown JH, Gillooly JF, Allen AP, Savage VM, West GB. 2004 \href{http://onlinelibrary.wiley.com/doi/10.1890/03-9000/full}{Toward a metabolic theory of ecology}. \emph{Ecology} \textbf{85}, 1771--1789.}

\leavevmode\vadjust pre{\hypertarget{ref-gouhier_synchrony_2010}{}}%
\CSLLeftMargin{45. }%
\CSLRightInline{Gouhier TC, Guichard F, Gonzalez A. 2010 Synchrony and {Stability} of {Food} {Webs} in {Metacommunities}. \emph{The American Naturalist} \textbf{175}, E16--E34. (doi:\href{https://doi.org/10.1086/649579}{10.1086/649579})}

\leavevmode\vadjust pre{\hypertarget{ref-williams_simple_2000}{}}%
\CSLLeftMargin{46. }%
\CSLRightInline{Williams RJ, Martinez ND. 2000 Simple rules yield complex food webs. \emph{Nature} \textbf{404}, 180--183. (doi:\href{https://doi.org/10.1038/35004572}{10.1038/35004572})}

\leavevmode\vadjust pre{\hypertarget{ref-christensen_ecopath_1992}{}}%
\CSLLeftMargin{47. }%
\CSLRightInline{Christensen V, Pauly D. 1992 {ECOPATH} {II} --- a software for balancing steady-state ecosystem models and calculating network characteristics. \emph{Ecological Modelling} \textbf{61}, 169--185. (doi:\href{https://doi.org/10.1016/0304-3800(92)90016-8}{10.1016/0304-3800(92)90016-8})}

\leavevmode\vadjust pre{\hypertarget{ref-loreau_biodiversity_2021}{}}%
\CSLLeftMargin{48. }%
\CSLRightInline{Loreau M \emph{et al.} 2021 Biodiversity as insurance: From concept to measurement and application. \emph{Biological Reviews} \textbf{96}, 2333--2354. (doi:\href{https://doi.org/10.1111/brv.12756}{10.1111/brv.12756})}

\leavevmode\vadjust pre{\hypertarget{ref-lefcheck_piecewisesem_2016}{}}%
\CSLLeftMargin{49. }%
\CSLRightInline{Lefcheck JS. 2016 {piecewiseSEM}: {Piecewise} structural equation modelling in r for ecology, evolution, and systematics. \emph{Methods in Ecology and Evolution} \textbf{7}, 573--579. (doi:\href{https://doi.org/10.1111/2041-210X.12512}{10.1111/2041-210X.12512})}

\leavevmode\vadjust pre{\hypertarget{ref-murphy_semeff_2021}{}}%
\CSLLeftMargin{50. }%
\CSLRightInline{Murphy M. 2021 \href{https://CRAN.R-project.org/package=semEff}{{semEff}: {Automatic} {Calculation} of {Effects} for {Piecewise} {Structural} {Equation} {Models}}. }

\leavevmode\vadjust pre{\hypertarget{ref-grace_structural_2006}{}}%
\CSLLeftMargin{51. }%
\CSLRightInline{Grace JB, Keeley JE. 2006 A {Structural} {Equation} {Model} {Analysis} {Of} {Postfire} {Plant} {Diversity} {In} {California} {Shrublands}. \emph{Ecological Applications} \textbf{16}, 503--514. (doi:\href{https://doi.org/10.1890/1051-0761(2006)016\%5B0503:ASEMAO\%5D2.0.CO;2}{10.1890/1051-0761(2006)016{[}0503:ASEMAO{]}2.0.CO;2})}

\leavevmode\vadjust pre{\hypertarget{ref-brooks_glmmtmb_2017}{}}%
\CSLLeftMargin{52. }%
\CSLRightInline{Brooks ME, Kristensen K, Benthem KJ van, Magnusson A, Berg CW, Nielsen A, Skaug HJ, Mächler M, Bolker BM. 2017 \href{https://journal.r-project.org/archive/2017/RJ-2017-066/index.html}{{glmmTMB} {Balances} {Speed} and {Flexibility} {Among} {Packages} for {Zero}-inflated {Generalized} {Linear} {Mixed} {Modeling}}. \emph{The R Journal} \textbf{9}, 378--400.}

\leavevmode\vadjust pre{\hypertarget{ref-vandermeer_coupled_2004}{}}%
\CSLLeftMargin{53. }%
\CSLRightInline{Vandermeer J. 2004 Coupled {Oscillations} in {Food} {Webs}: {Balancing} {Competition} and {Mutualism} in {Simple} {Ecological} {Models}. \emph{The American Naturalist} \textbf{163}, 857--867. (doi:\href{https://doi.org/10.1086/420776}{10.1086/420776})}

\leavevmode\vadjust pre{\hypertarget{ref-brose_predator_2019}{}}%
\CSLLeftMargin{54. }%
\CSLRightInline{Brose U \emph{et al.} 2019 Predator traits determine food-web architecture across ecosystems. \emph{Nature Ecology \& Evolution} \textbf{3}, 919--927. (doi:\href{https://doi.org/10.1038/s41559-019-0899-x}{10.1038/s41559-019-0899-x})}

\leavevmode\vadjust pre{\hypertarget{ref-gellner_consistent_2016}{}}%
\CSLLeftMargin{55. }%
\CSLRightInline{Gellner G, McCann KS. 2016 Consistent role of weak and strong interactions in high- and low-diversity trophic food webs. \emph{Nature Communications} \textbf{7}, 11180. (doi:\href{https://doi.org/10.1038/ncomms11180}{10.1038/ncomms11180})}

\leavevmode\vadjust pre{\hypertarget{ref-angelis_stability_1975}{}}%
\CSLLeftMargin{56. }%
\CSLRightInline{Angelis DLD. 1975 Stability and {Connectance} in {Food} {Web} {Models}. \emph{Ecology} \textbf{56}, 238--243. (doi:\href{https://doi.org/10.2307/1935318}{10.2307/1935318})}

\leavevmode\vadjust pre{\hypertarget{ref-gravel_trophic_2011}{}}%
\CSLLeftMargin{57. }%
\CSLRightInline{Gravel D, Massol F, Canard E, Mouillot D, Mouquet N. 2011 Trophic theory of island biogeography. \emph{Ecology Letters} \textbf{14}, 1010--1016. (doi:\href{https://doi.org/10.1111/j.1461-0248.2011.01667.x}{10.1111/j.1461-0248.2011.01667.x})}

\leavevmode\vadjust pre{\hypertarget{ref-barbier_generic_2018}{}}%
\CSLLeftMargin{58. }%
\CSLRightInline{Barbier M, Arnoldi J-F, Bunin G, Loreau M. 2018 Generic assembly patterns in complex ecological communities. \emph{Proceedings of the National Academy of Sciences of the United States of America} \textbf{115}, 2156--2161. (doi:\href{https://doi.org/10.1073/pnas.1710352115}{10.1073/pnas.1710352115})}

\leavevmode\vadjust pre{\hypertarget{ref-ghosh_temperature_2024}{}}%
\CSLLeftMargin{59. }%
\CSLRightInline{Ghosh S, Matthews B, Petchey OL. 2024 Temperature and biodiversity influence community stability differently in birds and fishes. \emph{Nature Ecology \& Evolution}, 1--12. (doi:\href{https://doi.org/10.1038/s41559-024-02493-7}{10.1038/s41559-024-02493-7})}

\leavevmode\vadjust pre{\hypertarget{ref-lee_ipcc_2023}{}}%
\CSLLeftMargin{60. }%
\CSLRightInline{Calvin K \emph{et al.} 2023 {IPCC}, 2023: {Climate} {Change} 2023: {Synthesis} {Report}. {Contribution} of {Working} {Groups} {I}, {II} and {III} to the {Sixth} {Assessment} {Report} of the {Intergovernmental} {Panel} on {Climate} {Change} {[}{Core} {Writing} {Team}, {H}. {Lee} and {J}. {Romero} (eds.){]}. {IPCC}, {Geneva}, {Switzerland}. (doi:\href{https://doi.org/10.59327/IPCC/AR6-9789291691647}{10.59327/IPCC/AR6-9789291691647})}

\leavevmode\vadjust pre{\hypertarget{ref-dornelas_assemblage_2014}{}}%
\CSLLeftMargin{61. }%
\CSLRightInline{Dornelas M, Gotelli NJ, McGill B, Shimadzu H, Moyes F, Sievers C, Magurran AE. 2014 Assemblage {Time} {Series} {Reveal} {Biodiversity} {Change} but {Not} {Systematic} {Loss}. \emph{Science} \textbf{344}, 296--299. (doi:\href{https://doi.org/10.1126/science.1248484}{10.1126/science.1248484})}

\leavevmode\vadjust pre{\hypertarget{ref-blowes_geography_2019}{}}%
\CSLLeftMargin{62. }%
\CSLRightInline{Blowes SA \emph{et al.} 2019 The geography of biodiversity change in marine and terrestrial assemblages. \emph{Science} \textbf{366}, 339--345. (doi:\href{https://doi.org/10.1126/science.aaw1620}{10.1126/science.aaw1620})}

\leavevmode\vadjust pre{\hypertarget{ref-danet_past_2024}{}}%
\CSLLeftMargin{63. }%
\CSLRightInline{Danet A, Giam X, Olden JD, Comte L. 2024 Past and recent anthropogenic pressures drive rapid changes in riverine fish communities. \emph{Nature Ecology \& Evolution}, 1--12. (doi:\href{https://doi.org/10.1038/s41559-023-02271-x}{10.1038/s41559-023-02271-x})}

\leavevmode\vadjust pre{\hypertarget{ref-olden_homogocene_2018}{}}%
\CSLLeftMargin{64. }%
\CSLRightInline{Olden JD, Comte L, Giam X. 2018 The {Homogocene}: A research prospectus for the study of biotic homogenisation. \emph{NeoBiota} \textbf{37}, 23--36. (doi:\href{https://doi.org/10.3897/neobiota.37.22552}{10.3897/neobiota.37.22552})}

\leavevmode\vadjust pre{\hypertarget{ref-wang_biotic_2021}{}}%
\CSLLeftMargin{65. }%
\CSLRightInline{Wang S \emph{et al.} 2021 Biotic homogenization destabilizes ecosystem functioning by decreasing spatial asynchrony. \emph{Ecology} \textbf{102}, e03332. (doi:\href{https://doi.org/10.1002/ecy.3332}{10.1002/ecy.3332})}

\end{CSLReferences}

\end{document}
